(1) 利用恒等式 \(e^{iz\cos u}=\sum_{m=-\infty}^\infty i^mJ_m(z)e^{imu}\)，其中 \(u=2\pi x/\Lambda\)，可写 \(t(x)=\sum_m c_me^{i2\pi mx/\Lambda}\)，系数 \(c_m=i^mJ_m(\varphi_0)\)。每个指数波数代表一个透射方向。严格的复相位因子 \(i^m\) 改变各级光波相位，单独求强度时才可取模而忽略。

(2) 光栅方程对正入射给 \(k\sin\theta_m=2\pi m/\Lambda\)，即 \(\sin\theta_m=m\lambda/\Lambda\)。在 \(\Lambda\gg\lambda\) 的近轴标量近似及能流投影近似相同条件下，各级效率 \(\eta_m=|c_m|^2=J_m^2(\varphi_0)\)。因为 \(J_{-m}=(-1)^mJ_m\)，有 \(\eta_{-m}=\eta_m\)，所以无左右偏置。若有级次接近掠射或变成倏逝波，则须引入矢量边界与通量因子，不能盲用此简化效率。

(3) 由于薄相位掩膜处处满足 \(|t(x)|^2=1\)，一周期 Parseval 恒等式给 \(\sum_m|c_m|^2=\Lambda^{-1}\int_0^\Lambda |t(x)|^2dx=1\)。因此 \(\sum_{m=-\infty}^\infty J_m^2(\varphi_0)=1\)。这体现理想相位光栅只重分配各衍射级能量，不发生吸收；有限口径只会把各离散级展宽为有限角宽，足够分离时积分能量关系仍保持。

(4) 当 \(|\varphi_0|\ll1\)，\(J_0(\varphi_0)=1-\varphi_0^2/4+O(\varphi_0^4)\)，\(J_1(\varphi_0)=\varphi_0/2+O(\varphi_0^3)\)。故零级 \(\eta_0=1-\varphi_0^2/2+O(\varphi_0^4)\)，一对一级各有 \(\eta_{\pm1}=\varphi_0^2/4+O(\varphi_0^4)\)；二级强度从 \(O(\varphi_0^4)\) 才开始。零级损失的 \(\varphi_0^2/2\) 正好被两个一级补足，提供能量级数的自洽检查。
